% This is text associated with the open textbook "Open Electromagnetics"
% See immediately below for author, date
% <License info goes here (boilerplate, not read by project compiler)>

%ORB TITLE Magnetic Field Intensity 
%ORB AUTHOR 0        # S.W. Ellingson, Virginia Tech (ellingson@vt.edu)
%ORB DATE 20180731   # yyyymmdd
%ORB ABSTRACT 
 
\label{m0012_Magnetic_Field_Intensity} % <-- Must match name of this file and module directory name.  Don't mess with this!
\index{magnetic field intensity|(}
\index{field!magnetic|(}
\index{field intensity, magnetic|see{magnetic field intensity}}

Magnetic field intensity ${\bf H}$ is an alternative description of the magnetic field in which the effect of material is factored out.  
For example, the magnetic flux density ${\bf B}$ (reminder: Section~\ref{m0003_Magnetic_Flux_Density}) due to a point charge $q$ moving at velocity ${\bf v}$ can be written in terms of the \emph{Biot-Savart Law}:
\index{Biot-Savart Law}
%
\begin{equation}
{\bf B} = \mu ~ \frac{q {\bf v}}{4\pi R^2} \times \hat{\bf R}
\label{m0012_eMatB}
\end{equation}
%
where $\hat{\bf R}$ is the unit vector pointing from the charged particle to the field point ${\bf r}$, $R$ is this distance, ``$\times$'' is the cross product, and
$\mu$ is the permeability of the material.
We can rewrite Equation~\ref{m0012_eMatB} as: 
%
\begin{equation}
{\bf B} \triangleq \mu {\bf H}
\label{m0012_eHDef}
\end{equation}
%
with:
%
\begin{equation}
{\bf H} =  \frac{q {\bf v}}{4\pi R^2} \times \hat{\bf R}
\label{m0012_eHHH}
\end{equation}
%
so ${\bf H}$ in homogeneous media 
\index{homogeneous (media)}
does not depend on $\mu$.

Dimensional analysis of Equation~\ref{m0012_eHHH} reveals that the units for ${\bf H}$ are amperes per meter (A/m).  
However, ${\bf H}$ does not represent surface current density,\footnote{The concept of current density is not essential to understand this section; however, a primer can be found in Section~\ref{m0101_Current_Distributions}.} as the units might suggest.  
While it is certainly true that a distribution of current (A) over some linear cross-section (m) can be described as a current density having units of A/m, ${\bf H}$ is associated with the magnetic field and not a particular current distribution.
Said differently, ${\bf H}$ can be viewed as a description of the magnetic field in terms of an equivalent (but not actual) current.

%ORB HIGHLIGHT
\begin{mdframed}[backgroundcolor=yellow!20,nobreak=true] 
The magnetic field intensity 
\index{magnetic field intensity!definition}
${\bf H}$ (A/m), defined using Equation~\ref{m0012_eHDef}, is a description of the magnetic field independent from material properties. 
\end{mdframed}

It may appear that ${\bf H}$ is redundant information given ${\bf B}$ and $\mu$, but this is true only in homogeneous
media. The concept of magnetic field intensity becomes important -- and decidedly not redundant -- when we encounter boundaries between media having different permeabilities.
As we shall see in Section~\ref{m0023_Boundary_Conditions_on_H}, boundary conditions on ${\bf H}$ constrain the component of the magnetic field which is tangent to the boundary separating two otherwise-homogeneous regions.
If one ignores the characteristics of the magnetic field represented by ${\bf H}$ and instead
considers only ${\bf B}$, then only the perpendicular component of the magnetic field is constrained.

The concept of magnetic field intensity also turns out to be useful in a certain problems in which $\mu$ is not a constant, but rather is a function of magnetic field strength.  In this case, the magnetic behavior of the material is said to be \emph{nonlinear}. 
\index{nonlinear (media)}
For more on this, see Section~\ref{m0058_Magnetic_Materials}. 
 
%%%%%%%%%%%%%%%%%%%

{\bf Additional Reading:}
\begin{itemize}
\item ``\href{https://en.wikipedia.org/wiki/Magnetic_field}{Magnetic field}'' on Wikipedia.
\item ``\href{https://en.wikipedia.org/wiki/Biot-Savart_law}{Biot-Savart law}'' on Wikipedia.
\end{itemize}

\index{magnetic field intensity|)}
\index{field!magnetic|)}
